$if(has-frontmatter)$
\frontmatter
$endif$

$if(title)$
\makeatletter
\twocolumn[
\maketitle
\begin{@twocolumnfalse}

\begin{abstract}
% Write your abstract here ----------------------------------------------------

\vspace{1em}

\normalsize \noindent We examine the economic tensions underlying content moderation decisions on social media platforms. Specifically, we investigate a fundamental paradox: toxic content may drive higher user engagement while simultaneously deterring advertisers concerned with brand safety. This tension is particularly consequential given the substantial negative externalities toxic content creates for society, including documented harms to mental health, increased polarization, and in extreme cases, offline violence. Analyzing over 53,000 headline variations from randomized controlled trials on Upworthy—a platform that created clickbait content specifically optimized for social media sharing—we find that headlines with higher likelihood of rejection by content moderators generate significantly greater user engagement, with each standard deviation increase in rejection likelihood corresponding to approximately 2 percent higher click-through rates. To further explain this "toxicity premium," we employ two complementary Digital In-Context Experiments. The first experiment tests whether varying levels of feed toxicity reveal a theorized non-monotonic relationship with user engagement. The second uses an incentive-compatible design to map advertisers' bidding responses to content toxicity. By studying these competing dynamics, we provide insights into the economic mechanisms driving platform content moderation decisions and contribute to the growing discourse on digital governance, demonstrating how platforms must navigate trade-offs between engagement optimization, societal welfare, and brand safety concerns.



% abstract completed  ----------------------------------------------------------
\end{abstract}

\vspace{3em}


\end{@twocolumnfalse}
]
\makeatother
$endif$